\begin{problem}[A one-point nonreduced local ring]
Let $R=k[\varepsilon]/(\varepsilon^n)$ with $n\ge2$.  Prove $R$ is local, compute its units, maximal ideal, residue field, and prime spectrum.  Show how the powers of the maximal ideal distinguish different $n$.
\end{problem}

\paragraph{Why this problem matters.}
Residue fields alone do not capture infinitesimal thickness; the entire local ring is needed.

\paragraph{Useful ideas before starting.}
Write every element as $a_0+a_1\varepsilon+\cdots+a_{n-1}\varepsilon^{n-1}$.

\paragraph{Guided hints.}
An element is a unit iff its constant term is nonzero.

\begin{solution}
If $a_0\neq0$, factor the element as $a_0(1+u)$ with $u$ nilpotent and invert $1+u$ by the finite geometric series.  If $a_0=0$, the element lies in $(\varepsilon)$ and is nilpotent, hence not a unit.  Thus the nonunits are exactly $(\varepsilon)$, so $R$ is local with residue field $R/(\varepsilon)=k$.  Every prime contains the nilpotent $\varepsilon$, and $(\varepsilon)$ is maximal, so the spectrum has one point.  The nilpotency index $\mathfrak m^n=0$ but $\mathfrak m^{n-1}\neq0$ distinguishes the rings.
\end{solution}

\paragraph{What happened mathematically.}
Same point, same residue field, different local algebra: this is the cleanest demonstration of infinitesimal structure.

\paragraph{Extensions.}
Compare lengths of the successive quotients $\mathfrak m^i/\mathfrak m^{i+1}$.

\paragraph{Provenance.}
New canonical synthesis problem based on the local-ring and residue-field corpus.
